\documentclass[../../../main.tex]{subfiles}

\begin{document}

\chapter{[A] Sequence and Series}

Sequence and series are one of the common topics in algebraic
combinatory and even just algebra problems. In this note, we review
some of the foundational topics and intuition on sequence and series,
then we discuss more advanced topics from recursion and different
type of sums.

\section{Prerequisites}

Problems involving sequence and series appeared a lot in competition
style problems. Let's review some of the topics and gain some
intuition on olympiad style sequence and series with Abel
transformation.

\subsection{Telescoping}

Telescoping is an essential skill that we learn when we first get
started with competition math. Because this content itself is
not too difficult, let's finish our review with an example problem.

\begin{problem}
{USAMTS Year 11 Round 4 Problem 3 \easy}
{USAMTS Year 11 Round 4 Problem 3}
Find $S$ defined as the following.
\[
S
= \sqrt{ 1 + \frac{1}{1^2} + \frac{1}{2^2} }
    + \sqrt{ 1 + \frac{1}{2^2} + \frac{1}{3^2} } + \cdots
    + \sqrt{ 1 + \frac{1}{1999^2} + \frac{1}{2000^2} }
\]
\end{problem}
\begin{solution}
First, let's rewrite the sum $S$ using the sum notation.
\[
S
= \sum_{i=1}^{1999}
    \sqrt{ 1 + \frac{1}{i^2} + \frac{1}{(i + 1)^2} }
\]
Now we see the terms $1 + \frac{1}{i^2} + \frac{1}{(i + 1)^2}$
inside a square root. To get rid of the annoying square root,
we can try writing the terms in a perfect square. Consider the
following equation.
\begin{align*}
    &\quad\ \left( 1 + \frac{1}{n} + \frac{1}{n + 1} \right)^2 \\
    &= 1 + \frac{1}{n^2} + \frac{1}{(n + 1)^2} + 2 \left(
        \frac{1}{n} + \frac{1}{n + 1} + \frac{1}{n(n + 1)}
    \right)
\end{align*}
Notice that because $\frac{1}{n} - \frac{1}{n+1} = \frac{1}{n(n+1)}$,
it would be much preferable to have $\frac{1}{n} - \frac{1}{n}
- \frac{1}{n(n+1)}$ instead of $\frac{1}{n} + \frac{1}{n + 1} +
\frac{1}{n(n + 1)}$. Therefore, the following equations are obtained.
\begin{align*}
    &\quad\ \left( 1 + \frac{1}{n} - \frac{1}{n + 1} \right)^2 \\
    &= 1 + \frac{1}{n^2} + \frac{1}{(n + 1)^2} + 2 \left(
        \frac{1}{n} - \frac{1}{n + 1} - \frac{1}{n(n + 1)}
    \right) \\
    &= 1 + \frac{1}{n^2} + \frac{1}{(n + 1)^2}
\end{align*}
Substituting the equation to the series, we obtain the following
equations.
\begin{align*}
    S
    &= \sum_{i=1}^{1999}
        \sqrt{1 + \frac{1}{i^2} + \frac{1}{(i + 1)^2}}
    = \sum_{i=1}^{1999} \sqrt{\left(
        \frac{1}{i} - \frac{1}{i + 1} + 1
    \right)^2} \\
    &= \sum_{i=1}^{1999}
        \left( \frac{1}{i} - \frac{1}{i + 1} + 1 \right)
    = 1999 + \frac{1}{1} - \frac{1}{2000} \\
    &= \frac{2000^2 - 1}{2000}
    = \frac{3999999}{2000}
\end{align*}
Thus, $S = \frac{3999999}{2000}$.
\end{solution}

I think telescoping when paired with inequalities like Cauchy
and products becomes much powerful! Continuing, let's discuss
Abel's transformation.

\subsection{Abel Transformation}

There are many different names of this, and when you search ``Abel
Transformation'' you would probably get different summation concerning
integration by parts. But let's stick with this term because I think
it is quite popular.

\begin{definition}{Abel Transformation}{Abel Transformation}
Given reals $a_i$ and $b_i$, the following equation holds where
$a_{n+1} = 0$.
\vspace{-0.6\baselineskip}
\[
\sum_{i=1}^n a_i b_i
= \sum_{i=1}^n (a_i - a_{i+1}) \left( \sum_{k=1}^i b_k \right)
\]
\end{definition}

This transformation is quite unique. To derive this, simply
expand the right-hand side. Then magically, every residual terms
telescope. To see how this can be applied in real life,
here is an example problem.

\begin{problem}{Classic Problem}{}
If $(a_1, \ldots, a_n) \succ (b_1, \ldots, b_n)$ with
$a_i \geq a_j \geq 0$ and $b_i \geq b_j \geq 0$ for $i > j$,
show that $\sum_{i=1}^n a_i^k \geq \sum_{i=1}^n b_i^k$ for all
positive integer $k$.
\end{problem}
\begin{solution}
We could use Karamata to prove this, but that would ruin all the fun!
For this problem, notice that we need to show that $\sum_{i=1}^n
(a_i^k - b_i^k) \geq 0$. Because the following equation
\[
a_i^k - b_i^k
= (a_i - b_i) \left(
    a_i^{k-1} + a_i^{k-2} b_i + \cdots + b_i^{k-1}
\right)
\]
holds, let $x_i = a_i^{k-1} + a_i^{k-2} b_i + \cdots + b_i^{k-1}$
and $y_i = a_i - b_i$. Rewriting $\sum_{i=1}^n (a_i^k - b_i^k)$,
it suffices to show that $\sum_{i=1}^n x_i y_i \geq 0$. By Abel
transform, we obtain following the equation.
\[
\sum_{i=1}^n x_i y_i
= \sum_{i=1}^n (x_i - x_{i+1}) (y_1 + \cdots + y_i)
\]
Notice that $x_i - x_{i+1} \geq 0$ since $a_i \geq a_j \geq 0$ and
$b_i \geq b_j \geq 0$ for $i > j$. Moreover, because
$(a_1, \ldots, a_n) \succ (b_1, \ldots, b_n)$, $y_1 + \cdots + y_i
> 0$. Therefore, the inequality holds.
\end{solution}

This is it for this section! In the next section, we will discuss
recurrence relation in depth.

\end{document}


